% CVPR introduction reframed after FlowAnchor (arXiv:2604.22586v1).
% Positioning: look-ahead demand estimation + normalized native-control allocation.
% METHODNAME is a placeholder. Add a quantitative headline only after the final
% official-FiVE and FiVE-C evaluations are frozen.

\section{Introduction}
\label{sec:intro}

Training-free video editing repurposes frozen video generators to transform a
source video according to a textual or visual instruction, without task-specific
fine-tuning or paired editing data~\cite{wan2025,selfforcing2025,flowedit,
tokenflow2024,streamedit2026}. Its central challenge is to realize the requested
change while preserving the source content and motion that the instruction does
not concern. Existing editors negotiate this trade-off through a native control:
for example, the step at which a FlowEdit-style velocity correction begins, or
the rate at which StreamEdit releases a target branch from source attention.
This control is ordinarily a single value shared by all locations and frames.

One global value is poorly matched to the local nature of an edit. Replacing a
rabbit with a drone requires some tokens to abandon the source silhouette early,
when structure forms; changing an object's colour requires less freedom; the
surrounding grass should remain tied to the source throughout. A strong global
setting may complete the replacement but disturb the background and source
motion, whereas a conservative one preserves the video but under-realizes the
edit. The problem is therefore not only to choose a better operating point on a
global editability--preservation curve, but to allocate editability according to
the different demands of individual video tokens.

Prior work provides two important forms of spatial adaptation. User-provided
masks and attention-derived supports constrain \emph{where} an edit acts
~\cite{objwiper,flowdirector}; Differential Diffusion assigns continuous
per-pixel strengths from an external change map~\cite{diffdiff}. Closest to our
setting, FlowAnchor~\cite{flowanchor} stabilizes inversion-free flow-based video
editing using a supplied spatial anchor: it refines cross-attention within that
region and, at each sampling step, amplifies the source--target velocity
difference according to its sample-wise normalized contrast and the video
length. This is an important instance of online, spatially adaptive signal
control. It remains tied, however, to the instantaneous velocity difference of
a FlowEdit trajectory. It does not produce a persistent estimate of local edit
demand before the final render, nor translate such an estimate into editors
whose native controls have different mechanisms and units.

We ask a complementary question: \textbf{can an editor look ahead at its own
behaviour, infer a continuous local demand for the requested transformation, and
then allocate its native edit budget accordingly?} Our key observation is that
a short, permissive edit can serve as a diagnostic intervention. When source
preservation is temporarily relaxed, the editor reveals the transformation it
is capable of making: where the output departs from the source provides an
output-space proxy for local edit demand. Unlike an online velocity magnitude,
this evidence is available before the final render, is defined for editors
without an additive editing signal, and can be placed on a fixed reference
scale rather than re-ranked independently within every video.

Based on this observation, we introduce \textbf{METHODNAME}, a look-ahead
controller for training-free video editing. METHODNAME first generates a short,
high-editability \emph{draft}. Source- and target-branch attention identifies
the content implicated by the instruction, requiring no user-provided or
benchmark mask. Within this support, semantic and perceptual draft--source
divergence estimates how strongly each token responded to the permissive edit.
We scale these distances by references frozen after calibration and shared
across videos and hosts, yielding a continuous edit-demand field without
per-video normalization. The draft is diagnostic only: its artifacts are not
carried into the final output.

Estimating demand and applying it are separate problems. A demand value cannot
be mapped linearly to arbitrary editor knobs: a query--key blend exponent and
an edit start step have different nonlinear effects over their respective
sampling schedules. We therefore derive, from each host's update rule, the mean
target-branch weight induced by every native control value. After normalizing
this response over the host's admissible control range, we invert it to assign a
native control to each token. Thus the same demand occupies the same normalized
target-exposure level in each host, while the host's editing mechanism remains
otherwise unchanged. The native-control mapping is analytic and is not fitted
to rendered quality or benchmark scores.

We instantiate METHODNAME on two substantially different editors.
StreamEdit~\cite{streamedit2026} is a few-step, visually prompted streaming
editor whose control governs a soft query--key blend between source and target
branches. Wan-Edit~\cite{five2025}, following FlowEdit~\cite{flowedit}, is
text-conditioned and controls a hard temporal gate on the source--target
velocity correction. This pairing tests whether one demand representation can
control both attention-based release and flow-based editing. We evaluate each
adaptive editor against its own global-control sweep at equal total sampling
budget, including the draft, and against support-only, region-level, and
draft-free alternatives. We report the official FiVE-Bench protocol and an
audited correction variant, FiVE-C, which preserves the videos, masks, edit
categories, and intended transformations while documenting repairs to textual
inconsistencies.
% TODO: Insert one quantitative headline supported by both matched-budget
% comparisons; do not claim an overall win if gains are limited to shape edits.

Our contributions are:
\begin{itemize}
    \item \textbf{Look-ahead edit-demand estimation.} We use a permissive draft
    as an output-space diagnostic and combine attention support with
    fixed-reference semantic and perceptual divergence to obtain a continuous,
    per-token demand field without an external mask.

    \item \textbf{Normalized native-control allocation.} We translate demand
    into heterogeneous editor controls by inverting a normalized
    target-exposure response derived from each host's update rule, without
    fitting the mapping to generated outputs.

    \item \textbf{Cross-mechanism, budget-matched evaluation.} We test the same
    demand estimator on attention-based StreamEdit and flow-based Wan-Edit,
    comparing each with tuned global and spatial-control alternatives while
    charging the diagnostic draft to the total sampling budget.
\end{itemize}
